\begin{figure*}[!t]
    \centering
    \begin{center}
    
    \begin{tcolorbox} [top=2pt,bottom=2pt, width=\linewidth, boxrule=1pt]
    {\footnotesize {\fontfamily{zi4}\selectfont
    \textbf{Object Selection Prompt:}
    You are an experienced room designer, please assist me in selecting large *floor*/*wall* objects and small objects on top of them to furnish the room. You need to select appropriate objects to satisfy the customer's requirements.
    You must provide a description and desired size for each object since I will use it to retrieve objects. If multiple identical items are to be placed in the room, please indicate the quantity and variance type (same or varied).
    Present your recommendations in JSON format:

    \{ 
        object\_name:\{ \\
            "description": a short sentence describing the object, \\
            "location": "floor" or "wall", \\
            "size": the desired size of the object, in the format of a list of three numbers, [length, width, height] in centimeters, \\
            "quantity": the number of objects (int), \\
            "variance\_type": "same" or "varied", \\
            "objects\_on\_top": a list of small children objects (can be empty) which are placed *on top of* this object. For each child object, you only need to provide the object\ name, quantity and variance type. For example, \{"object\_name": "book", "quantity": 2, "variance\_type": "varied"\}
        \}
    \} \\
    
    *ONE-SHOT EXAMPLE* \\
    
    Currently, the design in progress is *INPUT*, and we are working on the *ROOM\_TYPE* with the size of ROOM\_SIZE.
    Please also consider the following additional requirements: REQUIREMENTS.
    
    Here are some guidelines for you: \\
    1. Provide reasonable type/style/quantity of objects for each room based on the room size to make the room not too crowded or empty. \\
    2. Do not provide rug/mat, windows, doors, curtains, and ceiling objects which have been installed for each room. \\
    3. I want at least 10 types of large objects and more types of small objects on top of the large objects to make the room look more vivid. \\
    
    Please first use natural language to explain your high-level design strategy for *ROOM\_TYPE*, and then follow the desired JSON format *strictly* (do not add any additional text at the beginning or end).
    }
    \par}
    \end{tcolorbox}

    \begin{tcolorbox} [top=2pt,bottom=2pt, width=\linewidth, boxrule=1pt]
    {\footnotesize {\fontfamily{zi4}\selectfont
    \textbf{Layout Design Prompt:}
    You are an experienced room designer.
    Please help me arrange objects in the room by assigning constraints to each object.
    Here are the constraints and their definitions: \\
    1. global constraint: \\
        1) edge: at the edge of the room, close to the wall, most of the objects are placed here. \\
        2) middle: not close to the edge of the room.
    
    2. distance constraint: \\
        1) near, object: near to the other object, but with some distanbce, 50cm < distance < 150cm. \\
        2) far, object: far away from the other object, distance >= 150cm.

    3. position constraint: \\
        1) in front of, object: in front of another object. \\
        2) side of, object: on the side (left or right) of another object.
    
    4. alignment constraint:
        1) center aligned, object: align the center of the object with the center of another object.
    
    5. Rotation constraint:
        1) face to, object: face to the center of another object. \\
    
    For each object, you must have one global constraint and you can select various numbers of constraints and any combinations of them and the output format must be:
    object | global constraint | constraint 1 | constraint 2 | ...
    
    For example:
    sofa-0 | edge \\
    coffee table-0 | middle | near, sofa-0 | in front of, sofa-0 | center aligned, sofa-0 | face to, sofa-0 \\
    tv stand-0 | edge | far, coffee table-0 | in front of, coffee table-0 | center aligned, coffee table-0 | face to, coffee table-0 \\
    
    Here are some guidelines for you: \\
    1. I will use your guideline to arrange the objects *iteratively*, so please start with an anchor object which doesn't depend on the other objects (with only one global constraint). \\
    2. Place the larger objects first. \\
    3. The latter objects could only depend on the former objects. \\
    4. The objects of the *same type* are usually *aligned*. \\
    5. I prefer objects to be placed at the edge (the most important constraint) of the room if possible which makes the room look more spacious. \\
    6. Chairs must be placed near to the table/desk and face to the table/desk. \\
    Now I want you to design \{room\_type\} and the room size is \{room\_size\}. \\
    Here are the objects that I want to place in the \{room\_type\}:
    \{objects\} \\
    Please first use natural language to explain your high-level design strategy, and then follow the desired format *strictly* (do not add any additional text at the beginning or end) to provide the constraints for each object.
    }
    \par}
    \end{tcolorbox}
    
    \end{center}
    \caption{Prompt templates for Object Selection Module and Layout Design Module.}
    \label{fig:prompt-2}
    \vspace{-.3cm}
\end{figure*}



